\section{Arbitrary history-dependent tie-breaking in changing environments}

The next results apply to every finite class $H$ containing $f$, with
$h=\log|H|$. On each round the learner selects a consistent member of $H$
using the previous history, before seeing the next sampled input. It may
randomize and its tie rule may depend on all previous source identities,
inputs, labels, and predictions. Minimum-size consistent circuits satisfy
this requirement with
$H=\bigcup_{j\le s}\mathcal C_j$, provided the choice of a minimum circuit
does not depend on the current sampled input. In this case
\[
 h\le\log s+s\log16+2s\log(d+s+1).
\]
The input-independent choice condition matters: a different consistent
hypothesis chosen after seeing each input can have much greater aggregate
error than any individual consistent hypothesis.

\begin{theorem}[Adaptively interleaved random-order environments]\label{thm:random-pools}
Fix $K$ finite input pools of sizes $m_1,\ldots,m_K\ge1$; repeated inputs
inside or between pools are allowed. Independently uniformly permute each
pool. At each round a scheduler, using the full observed past but no
unrevealed entries, chooses a nonexhausted pool and reveals its next
entry. It may stop at any time. Then every learner just specified satisfies
\[
 \mathbb E M\le K+2(h+1)\sum_{i=1}^K H_{m_i-1},
 \qquad H_0=0.
\]
For $0<\delta<1$, put
\[
 a=h+\log\frac{2\sum_i m_i}{\delta},\qquad
 B_\delta=K+2a\sum_i H_{m_i-1}.
\]
With probability at least $1-\delta$,
\[
 M\le(e-1)B_\delta+\log(2/\delta).
\]
In particular, if each pool has at most $2^d$ entries, expectation is
$O(K(h+1)(d+1))$, and high probability adds only polynomial dependence on
$\log(K/\delta)$ and $d$.
\end{theorem}

\begin{proof}
For pool $i$, let $R_{i,n}$ be the maximum pool error fraction of a member
of $H$ agreeing with the first $n$ entries in that pool, ignoring all other
labels. A candidate with error fraction $r$ survives that prefix with
probability at most $(1-r)^n\le e^{-nr}$. Thus
\[
 \mathbb P(R_{i,n}>u)\le\min\{1,e^{h-nu}\},
 \qquad \mathbb E R_{i,n}\le(h+1)/n.
\]
At the next visit to pool $i$, the selected predictor is consistent with
that prefix; all of its pool errors lie in the $m_i-n$ remaining entries.
Its conditional error probability is at most
$m_iR_{i,n}/(m_i-n)$. Charging each used local count $(i,n)$ once and
enlarging to all possible local counts gives the expectation bound because
\[
 \sum_{n=1}^{m_i-1}\frac{m_i}{n(m_i-n)}=2H_{m_i-1}.
\]
The first visit to each pool costs at most one. This argument is unchanged
by arbitrary adaptive interleaving because each unused next entry remains
uniformly distributed among its own pool's remaining entries.

For high probability, a union bound makes all inequalities
$R_{i,n}\le a/n$, $1\le n<m_i$, hold except with probability $\delta/2$.
Stop the mistake process immediately before using a failed prefix. Until
that stopping time its total conditional error probabilities are at most
$B_\delta$. The exponential supermartingale used above bounds its mistakes
by $(e-1)B_\delta+\log(2/\delta)$ except with probability $\delta/2$.
On the simultaneous good-prefix event the process never stops, proving
the claim.
\end{proof}

\begin{theorem}[Interleaved iid environments]
Fix $K$ distributions $D_1,\ldots,D_K$. A scheduler uses the full past to
choose a source, then receives a fresh independent input from that source.
For every source having at least one nonzero disagreement risk in $H$, let
\[
 \varepsilon_i=
 \min\{D_i(g\ne f):g\in H,\ D_i(g\ne f)>0\}.
\]
Omit and reindex sources on which all hypotheses agree with $f$; $K$
now counts only the remaining sources. If there are none, $M_\infty=0$
and no probability budget is needed. Otherwise,
\[
 \mathbb E M_\infty\le
 \sum_{i=1}^K\{1+(h+1)(1+\log(1/\varepsilon_i))\}.
\]
For a high-probability bound set
\[
 J_i=1+\lceil\log_2(1/\varepsilon_i)\rceil,\quad
 L=\sum_iJ_i,\quad a=h+\log(2L/\delta),\quad
 B_\delta=(L+K)a+2K.
\]
Then, with probability at least $1-\delta$,
\[
 M_\infty\le(e-1)B_\delta+\log(2/\delta).
\]
On a finite support with smallest positive atom $p_i$, one may replace
$\varepsilon_i$ by $p_i$. In particular, sources uniform on arbitrary
subsets of $\{0,1\}^d$ give an expectation of $O(K(h+1)(d+1))$ even with
unlimited adaptive switching among disjoint regions.
\end{theorem}

The source need not be observed by the learner. In particular the theorem
covers arbitrary predictable mixture weights over the fixed $D_i$: choose
a latent source from the current weights and then draw from that source.
There is no bound on the number of switches or on the distance between
the source distributions.

\begin{proof}
Couple source $i$ to its own iid stream and let $R_{i,n}$ be the maximum
$D_i$-risk of a hypothesis surviving the first $n$ samples of that stream.
The same union bound gives
$\mathbb P(R_{i,n}>u)\le\min\{1,e^{h-nu}\}$.
The variable $R_{i,n}$ is either zero or at least $\varepsilon_i$.
For $u>0$, monotonicity of the integrand gives
\[
 \sum_{n\ge1}\min\{1,e^{h-nu}\}
 \le\int_0^\infty\min\{1,e^{h-tu}\}\,dt=(h+1)/u.
\]
Consequently Tonelli's theorem, using $\varepsilon_i$ in place of $u$
for $0<u<\varepsilon_i$, yields
\[
 \sum_{n\ge1}\mathbb E R_{i,n}
 \le(h+1)\left\{\int_0^{\varepsilon_i}\frac{du}{\varepsilon_i}
                   +\int_{\varepsilon_i}^1\frac{du}{u}\right\}
 =(h+1)(1+\log(1/\varepsilon_i)).
\]
The initial local round has risk at most one. On every visit to source $i$
after $n$ samples, the current consistent predictor has risk at most
$R_{i,n}$, regardless of observations from other sources. Each local count
is used at most once. Summing these bounds proves the expectation claim.

For high probability, for each $1\le j\le J_i$ put
$n_{i,j}=\lceil a2^j\rceil$. A union bound gives, simultaneously,
\[
 R_{i,n_{i,j}}\le2^{-j}
\]
except with probability $Le^{h-a}=\delta/2$. Since
$2^{-J_i}<\varepsilon_i$, the last inequality implies
$R_{i,n_{i,J_i}}=0$. By monotonicity of $R_{i,n}$, the initial block costs
at most $2a+1$ and each subsequent block costs at most $a+2^{-j}$, so
\[
 \sum_{n\ge0}R_{i,n}\le(J_i+1)a+2.
\]
Stop immediately before using a prefix which violates the applicable
block bound. Until stopping, the conditional error sum is at most
$B_\delta$. The same exponential-supermartingale bound and the good-event
union bound establish the stated infinite-horizon probability guarantee.
\end{proof}

\begin{remark}[Bounded arbitrary departures]
Either environment theorem remains valid with at most $A$ extra arbitrary
input rounds, with an additive $A$ in the actual mistake bound and its
expectation. The extra noiseless observations only shrink the version
space. The clean source counters and their latent orders are unchanged.
This includes a finite adversarial excursion followed by arbitrarily
many further source switches.
\end{remark}


\subsection{Randomizing within Hamming layers restores a polynomial bound}

For circuits over all two-input Boolean gates, or over the
AND/OR/NOT basis, a simple randomization of the input order changes
the same example. Let $f(x)=\bigwedge_{i=1}^d x_i$, present $1^d$
and $0^d$ first, and then present the remaining Hamming levels in
increasing order. Independently for each level, use a uniformly
random permutation of that level. The learner remains deterministic
and chooses a smallest consistent circuit using any rule depending
only on the preceding history.

For $k\ge2$, after all smaller levels have been labeled zero, the
globally minimum consistent functions are precisely the surviving
$k$-variable conjunctions. Here is the elementary reason. If $h$
depends on the variables in $R$, then $h(1_R)=h(1^d)=1$, so
$|R|\ge k$. A circuit with $r$ binary gates depends on at most
$r+1$ variables. Surviving $k$-conjunctions attain size $k-1$;
at that minimum size, $h$ depends on exactly $k$ variables.
Every non-all-ones assignment on these variables has a zero
extension of weight below $k$, so $h$ is their conjunction.
The same argument applies to AND/OR/NOT circuits with all gates
charged, since unary gates add no dependencies.

Suppose that $m$ weight-$k$ inputs remain. The learner's chosen
conjunction is wrong at precisely one of them, its own support
indicator. Conditional on the complete preceding history, the next
input is uniform over the $m$ remaining inputs. Its conditional
mistake probability is therefore exactly $1/m$, regardless of the
history-dependent deterministic tie rule. Consequently stage $k$
has exactly
\[
 H_{\binom dk}=\sum_{m=1}^{\binom dk}\frac1m
\]
expected mistakes. The first two inputs and the weight-one level
contribute between zero and $d+2$ mistakes, so
\[
 \sum_{k=2}^{d-1} H_{\binom dk}
 \;\le\; \mathbb E M
 \;\le\; d+2+\sum_{k=2}^{d-1} H_{\binom dk}
 \;=\;\Theta(d^2).
\]
For the upper bound, $H_N\le1+\ln N\le1+d\ln2$.
For the lower bound, for each integer
$d/4\le k\le d/2$,
$\binom dk\ge(d/k)^k\ge2^{d/4}$, hence
$H_{\binom dk}\ge(d/4)\ln2$; there are $\Omega(d)$
such levels. The assertion is asymptotic in $d$.

Thus the order still progresses strictly from low to high Hamming
weight after the initial positive example; only the order inside
each level is randomized. This suffices to reduce the deterministic
learner's expected mistake count from exponential to quadratic in
the target's minimum gate count $d-1$. The exact phase expectations
hold for every deterministic history-only tie rule. No analogous
upper bound is asserted here for pure NAND circuits, because the
NAND proof uses certificates rather than a classification of every
minimum consistent function.
